\documentclass[11pt, letterpaper]{article}
\input{../../homework.tex}
\usepackage{titlepageBU}
\title{Problem Set 1}
\courseID{PY 451}
\courseSection{A1}
\usepackage{listings}
\definecolor{codegreen}{rgb}{0,0.5,0}
\definecolor{codepurple}{rgb}{0.58,0,0.82}
\definecolor{backcolour}{rgb}{0.95,0.95,0.95}
\lstdefinestyle{h}{
    backgroundcolor=\color{backcolour},
    commentstyle=\color{codegreen},
    keywordstyle=\color{blue},
    numberstyle=\tiny\color{black},
    stringstyle=\color{codepurple},
    basicstyle=\ttfamily\footnotesize,
    breakatwhitespace=false,         
    breaklines=true,                 
    keepspaces=true,                 
    numbers=left,                    
    numbersep=5pt,                  
    showspaces=false,                
    showstringspaces=false,
    showtabs=false,
    tabsize=2
}
\lstset{style=h}
\begin{document}
\makeproblem
\section{The Zero Vector}
Let $V$ be a vector space over $\mathbb{C}$ with zero vector $\ket{0} $. Then
\[
	0 \ket{v} \overset{*}{=} (0+0) \ket{v}  \overset{4}{=} 0 \ket{v} +0 \ket{v} \overset{8,9}{\implies} 0 \ket{v} +(0 \ket{v} -0 \ket{v}) \overset{8}{=} 0 \ket{v} \overset{\star}{=} 0 \ket{v} - 0 \ket{v}  \overset{8}{=}\ket{0} 
.\]
Note that the step denoted by $*$ follows by the field axioms, rather than the vector space axioms. Additionally, the step denoted $\star$ follows by subtracting $ \ket{0}$ from both sides of the previous implication that $0 \ket{v} +0 \ket{v} =0 \ket{v} $; I just moved stuff around to make it fit on one line.
\section{Vector Spaces}
For all of these problems, let $V$ be the given set, and define addition and scalar multiplication pointwise; i.e. $(f+g)(x)=f(x)+g(x)$ and $(zf)(x)=z\cdot f(x)$ for all $f,g\in V$ and $z\in\mathbb{C}$.
\begin{enumerate}
	\item Yes.
		\begin{enumerate}
			\item Closure under addition: let $f,g\in V$. We have
				\[
					(f+g)(0)=f(0)+g(0)=0+0=0 \text{ and }(f+g)(L)=f(L)+g(L)=0+0=0
				.\]
				Therefore, $f+g\in V$;
			\item Closure under scalar multiplication: let $z\in\mathbb{C} $. We have
				\[(zf)(0)=z(f(0))=z0=0 \text{ and }  (zf)(L)=z(f(L))=z0=0,\]
 so $zf\in V$.
		\end{enumerate}
	\item No. Using the same definitions as before, note that for any $f\in V$, $f'(0)>0$. Note that $-1\in\mathbb{C}$. We have
		\[
			\frac{\mathrm{d}}{\mathrm{d}x} (-1f)(x)=-f'(x)
		.\]
		It follows that $(-1f)'(0)=-f'(0)<0$, so $-1f\notin V$.
	\item Yes. It suffices to show the only element of the space is the 0 function, because $0+0=0$ and $\alpha 0=0$ for all $0$ vectors in all vector spaces, and all scalars $\alpha $ in all fields. Let $f\in V$. Since $f$ has a vanishing first derivative everywhere, it follows that $f$ is continuous everywhere. Let $x\in [0,L]$ and suppose for contradiction that $f(x)>0$. Since $f$ is continuous on $[0,L]$, by the intermediate value theorem, $f$ takes on every value on the interval $[0, f(x)]$ at some point on the interval $[0,x]$. Since $f$ is continuous and $f(x) \gneq0$, we have that $f$ is increasing on some subinterval of $[0, f(x)]$; equivalently, $f$ has a positive first derivative on this interval. Therefore, $f(x)=0$ for all $x\in [0,L]$. Note that this implies $f(L/2)=0$, so that property is also satisfied. Therefore, $f\in V$ implies $f=0$, so $V=\left\{ 0 \right\} $ is the trivial vector space.
\end{enumerate}
\section{Eigenvalues}
\begin{enumerate}
	\item Notice that $M= \ket{u} \bra{u}$. Let $\lambda \ket{x} = M \ket{x} $. We have
		\[
			\lambda \ket{x} = \ket{u} \braket{u}{x} = \ket{u} \sum_{n=1}^{10} nx_n
		.\]
		We will use index notation to denote this as $\ket{u} nx_n $. We have
		\[
			\lambda x_k = u_k nx_n = knx_n
		.\]
		We have
		\[
			\lambda x_1 = nx_n
		.\]
		Additionally,
		\[
			\lambda x_k = knx_n = k(\lambda x_1) \implies x_k = kx_1
		.\]
		It follows that
		\[
			\lambda x_1 = n(nx_1)=x_1\sum_{n=1}^{10}n^2
		.\]
		Therefore,
		\[
			\lambda = \sum_{n=1}^{10} n^2 = 385
		.\]
		Moreover, note that the rank of the matrix is 1:
		\begin{lstlisting}[language=Python]		
#!/bin/python3

import numpy as np

u = np.arange(1, 11)
u = u.reshape(-1, 1)

print(np.linalg.matrix_rank(u @ u.T))
\end{lstlisting}
This program returns 1. The number of nonzero eigenvalues of a symmetric matrix is equal to its rank, and we implicitly assumed $\lambda \neq 0$ when we stated $x_k=kx_1$. If we suppose $\lambda =0$, then we have $nx_n=0$, since $k\neq 0$ for all $k$. We can define $x_k = (-1)^k1/k$, and note that
\[
	\sum_{n=1}^{10} (-1)^n\frac{n}{n}=5-5=0
.\]
Therefore, $0$ is also an eigenvalue for at least this singular vector.
\item I got a fever so I'm not going to do this the creative way and am just going to write a script.
	\begin{lstlisting}[language=Python]
#!/bin/python3

import numpy as np

u = np.arange(1, 11)
u = u.reshape(-1, 1)

v = np.arange(1, 11) / 2
v = v.reshape(-1, 1)

w = np.arange(1, 11)
w = (-1) ** w / w
w = w.reshape(-1, 1)

print(np.linalg.eigvals((u @ u.T) + (v @ v.T) + (w @ w.T)))
	\end{lstlisting}
	This outputs
	\begin{lstlisting}[language=Bash]
		[ 4.81250000e+02+0.00000000e+00j  1.54976773e+00+0.00000000e+00j
  7.78879923e-15+0.00000000e+00j -9.48602991e-15+3.10320984e-15j
 -9.48602991e-15-3.10320984e-15j  1.96268824e-15+1.23172745e-15j
  1.96268824e-15-1.23172745e-15j -2.11183098e-15+1.96936937e-15j
 -2.11183098e-15-1.96936937e-15j -6.23523255e-16+0.00000000e+00j]
	\end{lstlisting}
	Since this matrix is the sum of three rank 1 matrices, it's rank should be no more than 3, so it should have no more than 3 nonzero eigenvalues. This implies that the last 8 eigenvalues are floating point errors, since they are very small. Moreover, running the same code but changing the last line to print \verb|np.linalg.matrix_rank(...)|, we find that the rank of the matrix is precisely 2. Therefore, the only actual eigenvalues are
	\[
		481.25,\qquad 1.55
	.\]
\end{enumerate}
\end{document}
